% Page 023

\seite{23}

entsprechend mehr oder weniger geräumig und mit\ob
einem Hofraum umgeben. Hervorragende Gebäude,\ob
sowie schönere Häuser giebt es nicht. Die Ein\ohb
\margmark{Einwohner 212}%
wohnerzahl ist laut dieser Zählung 212.

In diesem Jahre wurde der im Jahre 1893 in Angriff\ob
genommene Neubau des Communications\ohb
weges zwischen Sefferweich und Seffern vollendet.

\margmark{die Gemarkung\\Sefferweich}%
Laut Gemeindeschätzung umfaßt die Gemarkung\ob
Sefferweich ca 1000 ha; dieselbe wird von den\ob
Fluren von Malbergweich, Nattenheim, Bickendorf,\ob
Seffern und Bierbach begrenzt. Die Verkehrswege\ob
sind in den letzten Jahren so verbessert worden, daß\ob
man die Verhältnisse als befriedigend bezeichnen\ob
kann. Anders liegen die Wasserverhältnisse.

\margmark{Der Dorf-Brunnen}%
Schon im Jahre 1891 und 92 suchte die Gemeinde\ob
mit einem Kostenaufwande von 4500 M.\ eine\ob
anscheinend billige Wasserleitung herzustellen, um\ob
der drückenden Wassernoth, der schon bei Bürger\ohb
versammlungen gedacht, zu steuern. Es zeigte sich jedoch\ob
recht bald, daß die bewilligte Summe unzureichend\ob
war und kaum mehr wie vor ist die Gemeinde auf\ob
Wasser angewiesen. In diesem Jahre mußte die allgemeine\ob
Noth mit fließendem Wasser zu den oben\ob
liegenden Brunnen\unsicher\ob
Wasser geholt wurde. Hierbei spielte sich eine kleine\ob
Scene ab, welche etwas heißblütige
\pend
